\section{Threats to Validity}
\label{sec:threats}

\paragraph{Serving model.} Each worker's slots advance at full speed independently, so total
throughput is the linear sum of per-slot rates. Real GPU concurrency does not work that way: when
we measured a real 4-way parallel serving process, per-request service time rose from
83 ms idle to 162 ms at sustained saturation ($1.95\times$).
In the saturated region our simulator therefore \emph{overestimates} throughput by roughly a
factor of two, which inflates the achievable throughput against a fixed SLO; the p95 comparisons
between policies are less affected because all policies share the error, but this is an
assumption, not a measurement.

\paragraph{Network model.} Transfers are serialized on a per-worker pipe and the transfer
\emph{rate} does not degrade as more requests share it; multiple concurrent transfers are FIFO
rather than processor-shared, and inter-worker links are independent (no shared uplink, no
incast, no fabric-level queueing, no congestion control). This \emph{underestimates} network cost
at high load, i.e.~makes our reported benefits conservative. The real-cluster testbed shapes
bandwidth in the application layer rather than with \texttt{tc}; the kernel path (socket buffers,
loopback of a container network) is real, the shaping is not.

\paragraph{Cluster evidence.} Our real deployment is a \emph{single} node: containers with CPU
quotas, no GPUs, no RDMA, no multi-host fabric. It can establish direction and can measure the
cost of reading state; it cannot establish production magnitudes. We quantify this rather than
assert it: repeat-to-repeat dispersion on an idle machine is 10--30 %
(\S\ref{sec:real}), so every real-cluster delta in this paper should be read with that error bar.

\paragraph{Workloads.} Each of the two traces represents one service class (Azure: the
\texttt{code} subset with very short outputs; BurstGPT: ChatGPT/GPT-4 conversational traffic with
strong bursts), and our ``workload classes'' are token-quantile slices rather than labels the
traces provide. Traces were downloaded truncated (the first 50\,MB of the Azure week, a 50\,MB
BurstGPT sample), which shortens coverage but not arrival-rate calibration, since rates are
re-fitted per simulated prefix.

\paragraph{Engineering estimates.} Four coefficients are declared estimates, not measurements:
prefill/decode work coefficients ($\alpha=1,\beta=4$), the payload model
$\kappa_c\in\{4\,\text{B},330\,\text{KB}\}$ per token, the injected staleness magnitudes, and the
SLO multiplier ($m=3$). We mitigate the first two by sweeping the payload model (which is the
dominant one) and by reporting network share, a quantity that does not depend on $\alpha,\beta$
in magnitude; the SLO enters only through miss-rate secondary metrics, and changes in $m$ cannot
affect the p95 analysis that carries our claims. Absolute values of \texttt{compute\_rate} are
not physically calibrated, so \emph{relative} comparisons between policies are meaningful while
absolute second-level latencies are not.

\paragraph{Coverage gaps we are aware of.} (i) The discounting rule of \S\ref{sec:discount} was
evaluated only at load intensity $\rho=0.6$; the intensity sweep we do have used the weaker $v_1$
score, whose sign is unstable across intensities ($-15.6\%$ to $+30.8\%$). Extending the
\emph{corrected} rule along the intensity axis is the largest gap in this paper.
(ii) At $N>4$ we cannot construct strongly aligned configurations, so the alignment effect is
measured over a narrowing range of $\rho$ (down to $0.24$ at $N=16$); the anti-aligned benefit
itself is stable, but the contrast narrows for reasons of experimental control rather than
mechanism. (iii) Observation error is modeled only as a constant lag plus an optional averaging
window; the real testbed shows that actual age is a random variable and that fetch cost depends
on load, neither of which is in the simulator.

\paragraph{Scope.} Nothing here should be read as a claim about a specific production
deployment. The contribution is a boundary: which configurations, payload regimes and freshness
levels make a network term worth including, and which make it worse than ignoring the network.

\paragraph{Real-cloud coverage.} The real LLM measurements use one 1.5B model family and A10
instances, with a fixed approximately 43 MB KV payload in B2/C. They cover two
cloud providers and one deployment path, not a representative sample of models, GPU generations,
regions, or production networks.

\paragraph{Unmeasured input-length sweep.} The planned D-mini experiment, which would have added
two input lengths to test whether the cross-cloud first-token cost scales with KV volume, was not
executed. The paper therefore uses simulation to study data-volume variation and treats the real
cloud results as mechanism and deployment anchors only; no real-cloud scaling law is claimed.

\paragraph{Descriptive cloud comparison.} B2 has four paired service rounds and C has two
directions, so their results are descriptive rather than population-level statistical estimates.
The $800$ B2 requests and $200$ C requests are not $1000$ independent experiments: requests are
nested inside service rounds or cloud directions. C also places the client at different clouds in
the two directions, and the measured bandwidth is asymmetric, so cross-cloud differences cannot
be attributed entirely to KV transfer.

\paragraph{Cloud heterogeneity.} Although both sides use A10 accelerators, the cloud hosts,
drivers, egress paths, tunnel behavior, and background conditions are not identical. The real
measurements therefore support feasibility and observed overhead under the stated configuration,
not a universal ranking of cloud providers or a production cost model.
